\section{Nietzsche’s Paradox}

There is a certain segment on the Internet who are, or claim to be, influenced by or followers of \textbf{Friedrich Nietzsche}, who then involve themselves in a paradox. Although Nietzsche claims that all opinions, conceptual schemes, or world views are merely perspectives, these bloggers nevertheless seem convinced that their particular perspective is absolutely true. Furthermore, as much as they decry the alleged absolutism of others, their efforts at evangelization and disdain for other perspectives reveal that same cast of mind.

\paragraph{Propaganda}


Now it may be that they are consciously publishing propaganda, even if it doesn't appear that way. In this case, we need to evaluate them by their effectiveness, for the purpose of propaganda is to promote a certain perspective, religion, or party platform. So what is their goal? Does the propaganda achieve that goal? Does it lead to an increase in power?

Propaganda has to be short, memorable, and emotion-laden. For example, we recently read a comment linking Christianity to democracy. This is powerful stuff, apparently, for neo-pagans who don't trouble to notice the ``democracy'' is a Greek concept predating Christianity, and historically Christian Europe opposed democracy until the advent of the so-called Enlightenment and French Revolution. Yet as propaganda, it is powerful because it matches his direct experience with contemporary neo-Christians and any intellectual correction would be long and tedious.

\paragraph{Scientific Training}


If not, then we must evaluate them as though they are attempts at the scientific (in the larger sense of the word) training of a certain intellectual elite. Any such perspective carries much more stringent requirements. The more of the following qualities that the perspective exhibits, the ``better'', that is, it demonstrates a stronger ``will to power''.


\begin{itemize}


\item It must ``save the appearances'', that is, it accounts for observable facts.

\item The more if ``falsifies'' opposing perspectives, the better. That is, it should demonstrate that other perspectives don't save all the appearances.

\item It is coherent, consistent, and complete.

\item It follows ``Occam's Razor'', that is, it doesn't bring in non-essential elements.

\item It accounts for more than the immediate data. In other words, it has application to domains other than the one in dispute.

\item It demonstrates a knowledge of history, culture, philosophy, and science.

\item It addresses objections --- known or possible --- in advance.

\item Its formulation is well written, following good grammar and rules of composition.

\item It is beautiful, inspiring, and succinct
\end{itemize}


\paragraph{Gnosis}


Of course, the Hermetist is aware of the techniques of propaganda and scientific indoctrination. They are restricted to the sub-lunar and stellar or astral realms, whereas the Hermetist communes with angels and approaches the True Sun. While he recognizes the value of science and employs perspectives to obtain desired results as necessary, he also recognizes a higher and transcendent knowledge which he calls Gnosis. He then attempts to express in in the terms of science.


\flrightit{Posted on 2011-06-14 by Cologero}

\begin{center}* * *\end{center}

\begin{footnotesize}
\begin{sffamily}
\texttt{Count Cagliostro on 2011-06-14 at 17:36 said:}

``While he recognizes the value of science and employs perspectives to obtain desired results as necessary, he also recognizes a higher and transcendent knowledge''

THAT'S IT!!!

\end{sffamily}
\end{footnotesize}

